\subsection*{参考文献}

{[}1{]} A.-L. CAUCHY, Œuvres, Paris (Gauthier-Villars), 1882-1932.

{[}2{]} N. H. ABEL, Œuvres, vol.~1, p.~223-224, edited by Sylow and Lie (Christiania) 1881.

{[}3{]} G. Ascoli, Sulle curve limiti di una varietà data di curve, Mem. Accad. Lincei (3), vol.~18 (1883), pp.~521-586.

{[}4{]} K. WEIERSTRASS, Mathematische Werke, Berlin (Mayer und Müller), 1894-1903:

\begin{enumerate}
\def\labelenumi{\alph{enumi})}
\item
  Zur Theorie der Potenzreihen, Bd. 1, pp.~67-74;
\item
  Über die analytische Darstellbarkeit sogenannter willkürlicher Functionen reeller Argumente, Bd. 3, pp.~1-37.
\end{enumerate}

{[}5{]} M. H. STONE, The generalized Weierstrass approximation theorem, Mathematics Magazine, 21 (1948), pp.~167-183 and 237-254.

{[}6{]} C. DE LA VALLÉE-POUSSIN, Leçons sur l'approximation des fonctions d'une variable réelle, Paris, (Gauthier-Villars), 1919.

引用数字依次表示章、节和小节或习题。

T : V, I, 2.

\(a^{x}\) （a 为任意实数 \textgreater0，x 为任意实数） : V, 4, I.

\(\log_{a}x\) （a、x 为实数 \textgreater0，\(a \neq i\) ） : V, 4, I.

\(R^{n}\) : VI, I, I.

\(d(x, y)\) （欧氏距离） : VI, 2, I.

\(||x||\) （欧氏范数） : VI, 2, I.

\((x|y)\) （内积） : VI, 2, 2.

\(S_{n}, B_{n}, R_{n}^{*}\) : VI, 2, 3.

\(P_{n}(R), P_{n}\) : VI, 3, I.

\(\tilde{R}\) : VI, 3, 3.

\(\infty\) （ \(\tilde{R}\) 的点） : VI, 3, 3.

\(P_{n,p}(R), P_{n,p}\) : VI, 3, 5.

\(r(G), d(G)\) （G 为 \(R^{n}\) 的任一闭子群） : VII, I, 2.

\(G^{*}\) （G 为 \(R^{n}\) 的任一子群） : VII, I, 3.

\(T^{n}\) : VII, I, 4.

C, i: VIII, I, I.

\(\mathcal{R}(z), \Im(z), \bar{z}, |z|\) （z 为复数）: VIII, I, I.

\(C^{*}, U\) : VIII, I, 3.

H : VIII, I, 4.

\(e(x) (= e^{2\pi ix})\) : VIII, 2, I.

\((\widehat{\Delta_{1}, \Delta_{2}})(\Delta_{1}, \Delta_{2}\text{ 射线})\) : VIII, 2, 2.

\(\delta, \varpi\) : VIII, 2, 2.

Am (z) （z 为复数）: VIII, 2, 2.

cos θ, sin θ, tan θ （θ 为角）: VIII, 2, 2.

\(\cos_{a}x\) , \(\sin_{a}x\) , \(\tan_{a}x\) , \(\cot_{a}x\) （x 为实数）: VIII, 2, 4.

cos x, sin x, tan x, cot x （x 为实数）: VIII, 2, 4.

\((\widehat{D_{1}, D_{2}})(D_{1}, D_{2}\text{ 直线})\) : VIII, 2, 6.

\(\delta_{0}\) : VIII, 2, 6.

\(C^{n}\) : VIII, 4, I.

\(P_{n}(C), \tilde{C}\) : VIII, 4, 3.

∞ （ \(\tilde{C}\) 的点） : VIII, 4, 3.

\[
\mathbf {P} _ {n, p} (\mathbf {C}): \text { VIII }, 4, 4.
\]

\(\beta \mathbf{X}\) （X 为完全正则空间）: IX, 1, 习题 7.

\(d(\mathbf{A}, \mathbf{B}), d(x, \mathbf{A})\) {[} \(x\) 为一点，A、B 为一个度量空间的子集，其中两点间的距离记作 \(d(x, y)] : \mathrm{IX}, 2, 2\)

\(|x|\) （赋值除环中的绝对值）: IX, 3, 2.

\textbar\textbar x\textbar\textbar{} （赋范空间中的范数） : IX, 3, 3.

D(A) : IX, 5, 习题 3.

L(C) （C 为筛）: IX, 6, 5.

\(\prod_{n \in \mathbb{N}} x_n : \text{IX, Appendix, I.}\)

\(P_{n=h} x_n : \text{IX, Appendix, 3.}\)

\(\prod_{n=0}^{n} x_n: \text{IX, Appendix, 3.}\)

\(\mathcal{F}(\mathrm{X};\mathrm{Y}),\mathrm{H}(x)\) {[}H 为 \(\mathcal{F}(\mathrm{X};\mathrm{Y})]\) 的子集，\(\Phi (x)\) {[}\(\Phi\) 为 \(\mathcal{F}(\mathrm{X};\mathrm{Y})]\) 上的一个滤子，\(u|\mathrm{A},\mathrm{H}|\mathrm{A}\) {[}A 为 \(\mathrm{X},u\in \mathcal{F}(\mathrm{X};\mathrm{Y}),\mathrm{H}\subset \mathcal{F}(\mathrm{X};\mathrm{Y})]:\mathrm{X},\) 1.

\(W(\mathbf{V})\) （V 为 \(\mathbf{Y}\) 的一个围域） : \(\mathbf{X}, \mathrm{I}, \mathrm{I}\) .

\[
\mathcal {F} _ {u} (\mathrm{X}; \mathrm{Y}): \mathrm{X}, \mathrm{I}, \mathrm{I}.
\]

\(\mathcal{F}_{\mathfrak{S}}(\mathbf{X};\mathbf{Y}):\mathbf{X},\mathbf{1},\mathbf{2}.\)

\(W(A, V) (A \in \mathfrak{S}, V \text{ 是一个围域，相对于 } Y): X, I, 2.\)

\[
\mathcal {F} _ {s} (\mathrm{X}; \mathrm{Y}), \mathcal {F} _ {c} (\mathrm{X}; \mathrm{Y}): \mathrm{X}, \mathrm{I}, 3.
\]

\[
\mathcal {C} (\mathbf {X}; \mathbf {Y}), \mathcal {C} _ {\mathfrak {S}} (\mathbf {X}; \mathbf {Y}), \mathcal {C} _ {s} (\mathbf {X}; \mathbf {Y}), \mathcal {C} _ {c} (\mathbf {X}; \mathbf {Y}), \mathcal {C} _ {u} (\mathbf {X}; \mathbf {Y}): \mathbf {X}, \mathrm{I}, 6.
\]

\[
\tilde {\mathcal {C}} _ {\mathfrak {S}} (\mathrm{X}; \mathrm{Y}): \mathrm{X}, \mathrm{I}, 6.
\]

\(\tilde{x}\) （x 为 X 的点）: X, 2, 1.

\(\mathcal{B}_{\mathfrak{S}}(\mathbf{X};\mathbf{Y}),\mathcal{B}(\mathbf{X};\mathbf{Y})\) （Y 为度量空间）: \(\mathbf{X},3,\mathrm{i}\)

\(||u||\) （u 为到赋范空间中的有界映射）: X, 3, 2.

\(\mathscr{L}\left(\mathbf{X}_{1},\ldots,\mathbf{X}_{n};\mathbf{Y}\right)\left(\mathbf{X}_{1},\ldots,\mathbf{X}_{n}\text{ 和 Y 均为赋范空间}\right):\mathbf{X},3,2.\)

\(||u||\) （u 为赋范空间的乘积到赋范空间中的多重线性映射）: X, 3, 2.

\(T(\mathbf{K},\mathbf{U})\) （K 为 \(\mathbf{X}\) 的紧子集，U 为 \(\mathbf{Y}\) 的开子集） : \(\mathbf{X}\) , 3, 4. \(\mathcal{C}_{\beta}:\mathbf{X}\) , 3, 5.

\(\mathcal{C}^{\infty}\) (X; R): X, 4, 习题 6.

\(\nu X:X,4,\) 习题 17.

引用数字依次表示章、节和小节或习题。

横坐标 : VI, 1, 4. 横轴 : VI, 1, 4. 绝对值 : IX, 3, 2. 平凡绝对值 : IX, 3, 2. p 进绝对值 : IX, 3, 2. 等价绝对值 : IX, 3, 2. 复数的绝对值 : VIII, 1, 1. 复数的绝对收敛无穷乘积 : VIII, 3, 3. 绝对收敛级数（在赋范空间中）: IX, 3, 6. R\^{}n 中点的绝对收敛级数 : VII, 3, 2. 绝对可和族（在赋范空间中）: IX, 3, 6. 锐角域 : VIII, 2, 5. R\^{}n 的加法群 : VI, 1, 2. 模 a 的实数加法群 : V, 1, 2. R\^{}n 的加法一致结构 : VII, 1, 2. 容许子集 : IX, 6, 习题 12. 相互垂直的仿射线性簇 : VI, 2, 2. 赋范代数 : IX, 3, 7. 复数的代数范数 : VIII, 1, 1. 代数数 : VIII, 1, 习题 2. 几乎开映射 : IX, 5, 习题 24. 几乎开集 : IX, 5, 习题 6. 概周期函数 : X, 3, 习题 28. 概周期点 : X, 3, 习题 27. 复数的辐角 : VIII, 2, 2 与 VIII, 2, 3. 平角 : VIII, 2, 2. 一对射线的角 : VIII, 2, 2 与 VIII, 2, 3. 角域的角 : VIII, 2, 5. 正直角 : VIII, 2, 2.

角域（锐的、直的、钝的、平的、凸的、凹的）: VIII, 2, 5. 一致逼近: X, 4, 1. 阿基米德的（线性有序群）: V, 3, 习题 1. 阿基米德公理: V, 2. 复数的幅角: VIII, 2, 2. 阿斯科利定理: X, 2, 5. \(\mathbf{R}^n\) 的子群的相伴子群 : VII, 1, 3. 阿基米德公理: V, 2. 坐标轴: VI, 1, 4. 虚轴: VIII, 1, 2. 横轴: VI, 1, 4. 纵轴: VI, 1, 4. 实轴: VIII, 1, 2.

紧收敛的一致结构 : X, 1, 3.

紧开拓扑 : X, 3, 4.

斯通–切赫紧化 : IX, 1, 习题 7.

相容的（度量与拓扑） : IX, 2, 5.

相容的（度量与一致结构） : IX, 2, 4.

相容的（范数与代数结构） : IX, 3, 7.

完全正规空间 : IX, 4, 习题 3.

完全正则滤子 : IX, 1, 习题 8.

完全正则空间 : IX, 1, 5.

与可一致化空间相伴的完全正则空间 : IX, 1, 习题 4.

完全分离的（完全正则空间中的闭集）: IX, 1, 习题 11.

\(\mathbf{C}^n\) 中的复超平面 : VIII, 4, 1.

\(\mathbf{C}^n\) 中的复线性簇 : VIII, 4, 1.

\(\mathbf{C}^n\) 中的复直线 : VIII, 4, 1.

复数 : VIII, 1, 1.

n 维复数空间 : VIII, 4, 1.

\(\mathbf{C}^n\) 中的复平面 : VIII, 4, 1.

复射影直线 : VIII, 4, 3.

复射影平面 : VIII, 4, 3.

n 维复射影空间 : VIII, 4, 3.

复数的共轭 : VIII, 1, 1.

连续单位分解 : IX, 4, 3.

收敛的无穷乘积（在赋范代数中） : IX, 附录, 3.

正规收敛的 : X, 3, 2.

逐点收敛的 : X, 1, 3.

一致收敛的 : X, 1, 1.

坐标轴 : VI, 1, 4.

坐标簇 : VI, 1, 4.

角的余弦 : VIII, 2, 2.

数的余弦 : VIII, 2, 4.

角的余切 : VIII, 2, 2.

数的余切 : VIII, 2, 4.

可数型（可度量化空间） : IX, 2, 8.

可数紧空间 : IX, 2, 习题 14.

可除覆盖 : IX, 4, 习题 16.

均匀覆盖 : IX, 4, 习题 16.

逐点有限覆盖 : IX, 4, 3.

一对直线的叉角 : VIII, 2, 6.

直叉角 : VIII, 2, 6.

方体 : IX, 1, 5.

立方体（开的、闭的） : VI, I, I.

佩亚诺曲线 : VI, 1, 习题 2.

度（角度量的单位） : VIII, 2, 3.

直径 : IX, 2, 2.

直径超平面 : VI, 2, 4.

\(\mathbf{S}_n:\mathbf{VI},3,\mathbf{i}\) 的对径点

\(\mathbf{R}^n\) 的闭子群的维数 : VII, 1, 2.

迪尼定理 : X, 4, 1.

直线或射线的方向比 : VI, 1, 4.

直线或射线的方向向量 : VI, 1, 4.

圆盘（开的、闭的） : VI, 2, 3.

子集的离散族 : IX, 4, 习题 18.

欧氏位移 : VI, 2, 2.

两集合间的距离 : IX, 2, 2.

欧氏距离 : VI, 2, 1.

点到集合的距离 : IX, 2, 2.

可除覆盖 : IX, 4, 习题 16.

四元数体 : VIII, 1, 4.

赋值除环 : IX, 3, 2.

\(\varepsilon\) -周期（函数的）: X, 3, 习题 28.

等度连续的映射集 : X, 2, 1.

在一点等度连续 : X, 2, 1.

一致等度连续 : X, 2, 1.

模 i 均分 : VII, i, 习题 14.

等价绝对值 : IX, 3, 2.

等价的伪度量族 : IX, 1, 2.

等价范数 : IX, 3, 3.

等价伪度量 : IX, 1, 2.

等价半绝对值 : IX, 3, 习题 10.

映入 \(\mathbf{P}_n\) 的本质映射 : VI, 3, 习题 2.

映入 \(\mathbf{S}_1\) 的本质映射 : VI, 2, 习题 6.

球（闭的、开的） : VI, 2, 3.

欧氏位移 : VI, 2, 2.

欧氏距离 : VI, 2, 1.

\(\mathbf{R}^n\) 中的欧氏范数 : VI, 2, 1.

球面 : VI, 2, 3.

均匀覆盖 : IX, 4, 习题 16.

指数函数 : V, 4, 1.

乘积的因子（在赋范代数中） : IX, 附录, 1.

绝对可和族 : IX, 3, 6.

从属于子集族的函数族 : IX, 4, 3.

法里数列 : VII, 1, 习题 13.

复数域 : VIII, 1, 1.

完全正则滤子 : IX, 1, 习题 8.

极大完全正则滤子 : IX, 1, 习题 8.

有限开覆盖的一致结构 : IX, 4, 习题 17.

平角 : VIII, 2, 2.

平角域 : VIII, 2, 5.

概周期函数 : IX, 3, 习题 28.

指数函数 : V, 4, 1.

周期函数（定义在 \(\mathbf{R}^n\) 上）: VII, 1, 6.

q 重周期函数 : VII, 1, 6.

由一组子集生成的（σ 代数） : IX, 6, 3.

百分度（角度量的单位） : VIII, 2, 3.

加法群（ \(\mathbf{R}^n\) 的）: VI, 1, 2.

加法群（模 a 的实数的）: V, 1, 2.

可度量化群 : IX, 3, 1.

欧氏位移群 : VI, 2, 2.

单参数群 : V, 3.

正交群 : VI, 2, 2.

半直线 : VI, 1, 4.

由超平面确定的半空间（开的、闭的） : VI, 1, 4.

豪斯多夫半绝对值 : IX, 3, 习题 9.

与半绝对值相伴的豪斯多夫半绝对值 :

半球面（开的、闭的） : VI, 2, 4.

分式线性函数 : VI, 3, 5.

超实极大理想 : X, 4, 习题 16.

无穷远超平面 : VI, 3, 3.

复超平面（在 \(\mathbf{C}^n\) 中）: VIII, 4, 1.

直径超平面 : VI, 2, 4.

投影超平面（在球极投影中） : VI, 2, 4.

虚轴 : VIII, 1, 2.

\(\mathbf{R}^n\) 中的格 : VII, I, I.

左不变度量 : IX, 3, 1.

左不变伪度量 : IX, 3, 习题 1.

\(\mathbf{R}^n:\mathrm{VI},2,\mathrm{i}\) 中线段的长度

折线 : VI, 1, 习题 6.

简单折线 : VI, 1, 习题 9.

复射影直线 : VIII, 4, 3.

实射影直线 : VI, 3, 1.

函数的线性组合（系数在赋范空间中）: X, 4, 4.

复线性簇（在 \(\mathbf{C}^n\) 中）: VIII, 4, 1.

线性可达的 : IX, 6, 习题 11.

复直线（在 \(\mathbf{C}^n\) 中）: VIII, 4, 1.

局部仿紧空间 : IX, 4, 习题 27.

以 \(a:\mathbf{V},4,\mathrm{i}\) 为底的对数

对数螺线 : VIII, 3, 习题 5.

卢津空间 : IX, 6, 4.

几乎开映射 : IX, 5, 习题 24.

波莱尔映射 : IX, 6, 习题 16.

由筛分诱导的映射 : IX, 6, 5.

映入 \(\mathbf{P}_n\) 的映射，本质的或非本质的 : VI, 3, 习题 2.

映入 \(S_{1}\) 的映射，本质的或非本质的 : VI, 2, 习题 6.

等距映射 : IX, 2, 2.

分段线性映射 : VI, 1, 习题 6.

极大完全正则滤子 : X, 1, 习题 8.

贫集 : IX, 5, 2.

角的测度 : VIII, 2, 3.

叉角的测度 : VIII, 2, 6.

主测度（角的）: VIII, 2, 3.

主测度（叉角的）: VIII, 2, 6.

亚紧空间 : IX, 4, 习题 25.

度量 : IX, 2, 1.

与伪度量相伴的度量 : IX, 2, 1.

与拓扑相容的度量 : IX, 2, 5.

与一致结构相容的度量 : IX, 2, 4.

左（右）不变度量 : IX, 3, 1.

度量空间 : IX, 2, 1.

可度量化拓扑群 : IX, 3, 1.

可度量化拓扑空间 : IX, 2, 5.

可数型的可度量化拓扑空间 : IX, 2, 8.

可度量化一致空间 : IX, 2, 4.

可度量化一致结构: IX, 2, 4.

可乘序列（在赋范代数中） : IX, 附录, 1.

长田–斯米尔诺夫定理 : IX, 4, 习题 22.

负实半轴 : VIII, 1, 2.

非贫空间 : IX, 5, 习题 7.

范数（在向量空间上）: IX, 3, 3.

代数范数（复数的）: VIII, 1, 1.

与代数结构相容的范数 : IX, 3, 7.

欧氏范数（在 \(\mathbf{R}^n\) 中）: VI, 2, 1.

正规空间 : IX, 4, 1.

正规收敛级数 : X, 3, 2.

赋范代数 : IX, 3, 7.

赋范空间 : IX, 3, 3.

等价范数 : IX, 3, 3.

无处稠密集 : IX, 5, 1.

代数数 : VIII, 1. 习题 2.

复数 : VIII, 1, 1.

n 维复数空间 : VIII, 4, 1.

n 维实数空间 : VI, I, I.

钝角域 : VI, 2, 5.

单参数群 : V, 3.

开球 : IX, 2, 2.

反向射线 : VI, 1, 4.

有序分划 : IX, 附录, 习题 2.

纵坐标 : VI, I, 4.

纵轴 : VI, 1, 4.

射线的端点 : VI, 1, 4.

相互垂直的仿射线性簇 : VI, 2, 2.

正交群 : VI, 2, 2.

正交变换 : VI, 2, 2.

正交向量、正交向量子空间 : VI, 2, 2.

函数在一点的振幅 : IX, 2, 3.

函数在集合中的振幅 : IX, 2, 3.

\(p\) -进绝对值: IX, 3, 7.

超平行体（开的、闭的） : VI, 1, 3.

连续单位分解 : IX, 4, 3.

有序分划 : IX, 附录, 习题 2.

佩亚诺曲线 : VI, 1, 习题 2.

完备正规空间 : IX, 4, 习题 7.

周期函数（定义在 \(\mathbf{R}^n\) 上）: VII, 1, 6.

周期函数的周期 : VII, 1, 6.

分段线性映射 : VI, I, 习题 6.

抽屉原理 : VII, 1, 习题 11.

复射影平面 : VIII, 4, 3.

实射影平面 : VI, 3, 1. 复平面（在 C\(^{n}\) 中） : VIII, 4, 1. 无穷远点 : VI, 3, 3. 逐点有限覆盖 : IX, 4, 3. 一致收敛点 : X, 1, 习题 9. 在 X 的子集上的逐点收敛 : X, 1, 3. 逐点收敛的拓扑 : X, 1, 3 与 3, 4. 逐点收敛的一致结构 : X, 1, 3. 逐点收敛的（滤子） : X, 1, 3. 波兰空间 : IX, 6, 1. 属于给定集合的函数的多项式 : X, 4, 2 与 4, 4. 三角多项式 : X, 4, 4. 正实半轴 : VIII, 1, 2. 预紧收敛的一致结构 : X, 1, 3. 角的主测度 : VIII, 2, 3. 叉角的主测度 : VIII, 2, 6. q 重周期函数的基本周期系 : VII, 1, 6. 绝对收敛的乘积（复数的） : VIII, 3, 3. 无穷乘积 : IX, 附录, 3. 赋范代数中可乘序列的乘积 : IX, 附录, 1. 内积 : VI, 2, 2. 中心投影 : VI, 2, 3. 投影超平面 : VI, 2, 4. 球极投影 : VI, 2, 4. 投影顶点 : VI, 2, 4. n 维复射影空间 : VIII, 4, 3. n 维实射影空间 : VI, 3, 1. 在一点真作用（算子群） : X, 2, 习题 18. 伪紧空间 : IX, 1, 习题 22. 伪度量 : IX, 1, 1. 不变伪度量 : IX, 3, 习题 1. 等价伪度量 : IX, 1, 2. 纯虚数（复数） : VIII, 1, 1.

R\(^{n}\) 上的 q 重周期函数 : VII, 1, 6. 二次超曲面（在 P\(_{n}\) 中）: VI, 3, 习题 10. 二次曲面（在 R\(^{n}\) 中）: VI, 2, 习题 10. 二次锥面（在 P\(_{n}\) 中）: VI, 3, 习题 11. 二次锥面（在 R\(^{n}\) 中）: VI, 2, 习题 11. 四元数 : VIII, 1, 4.

弧度测度 : VIII, 2, 3. 球的半径 : IX, 2, 2. 球或球面的半径 : VI, 2, 3 与 IX, 2, 2.

\(\mathbf{R}^n\) 的子群的有理秩 : VII, 1.

方向比 : VI, 1, 4.

射线（闭的、开的） : VI, 1, 4.

反向射线 : VI, 1, 4.

实轴 : VIII, 1, 2.

实紧空间 : X, 4, 习题 17.

空间的实紧化 : X, 4, 习题 17.

实线性簇（在 \(\mathbf{C}^n\) 中）: VIII, 4, i.

实极大理想: X, 4, 习题 16.

\(n\) 维实数空间 : VI, I, I.

复数的实部 : VIII, 1, 1.

实射影直线 : VI, 3, 1.

实射影平面 : VI, 3, 1.

\(n\) 维实射影空间 : VI, 3, 1.

除环的实赋值 : IX, 3, 2.

矩形（开的、闭的） : VI, I, I.

凹角域 : VIII, 2, 5.

正直角 : VIII, 2, 2.

直角域 : VIII, 2, 5.

直叉角 : VIII, 2, 6.

右不变度量 : IX, 3, 1.

右不变伪度量 : IX, 3, 习题 1.

\(\sigma\) -代数: IX, 6, 3.

凸角域 : VIII, 2, 5.

伪度量的饱和族 : IX, 1, 2.

子集的饱和集 : X, 1, 习题 5.

内积 : VI, 2, 2.

角域 : VIII, 2, 5.

线段（闭的、开的、在 x 处开而在 y 处闭的） : VI, 1, 4.

线段的长度 : VI, 2, 1.

半绝对值 : IX, 3, 习题 9.

豪斯多夫半绝对值 : IX, 3, 习题 9.

等价半绝对值 : IX, 3, 习题 10.

负实半轴 : VIII, 1, 2.

正实半轴 : VIII, 1, 2.

半圆（开的、闭的） : VI, 2, 4.

可分可度量化空间 : IX, 2, 8.

分离的函数集 : X, 4, 1.

集合的点被函数集分离 : X, 4, 1.

可乘序列 : IX, 附录, 1.

绝对收敛级数（ \(\mathbf{R}^n\) 的点的）: VII, 3, 2.

法里数列 : VII, 1, 习题 13.

几乎开集 : IX, 5, 习题 6.

波莱尔集 : IX, 6, 3.

有界集 : IX, 2, 2.

有界集（在 \(\mathbf{R}^n\) 中）: VI, I, I.

可容集 : IX, 6, 9.

贫集 : IX, 5, 2.

无处稠密集 : IX, 5, 1.

苏斯林集 : IX, 6, 2.

星形集（在 \(\mathbf{R}^n\) 中）: VI, 2, 习题 12.

瘦集 : IX, 5, 习题 2.

星形集的壳 : VI, 2, 习题 12.

立方体的面 : VI, I, I.

折线的边 : VI, 1, 习题 6.

筛 : IX, 6, 5.

筛分 : IX, 6, 5.

由筛分诱导的映射 : IX, 6, 5.

简单折线 : VI, 1, 习题 9.

角的正弦 : VIII, 2, 2.

数的正弦 : VIII, 2, 4.

截片（在线性有序集中）: IX, 附录, 习题 2.

直线的斜率 : VIII, 2, 6.

苏斯林集 : IX, 6, 2.

苏斯林空间 : IX, 6, 2.

与筛分相伴的空间 : IX, 6, 5.

贝尔空间 : IX, 5, 3.

集体正规空间 : IX, 4, 习题 18.

完全正规空间 : IX, 4, 习题 3.

完全正则空间 : IX, 1, 5.

复数空间 : VIII, 4, 1.

复射影空间 : VIII, 4, 3.

可数紧空间 : IX, 2, 习题 14.

局部仿紧空间 : IX, 4, 习题 27.

卢津空间 : IX, 6, 4.

亚紧空间 : IX, 4, 习题 25.

度量空间 : IX, 2, 1.

可数型的可度量化空间 : IX, 2, 8.

可度量化拓扑空间 : IX, 2, 5.

可度量化一致空间 : IX, 2, 4.

非贫空间 : IX, 5, 习题 7.

正规空间 : IX, 4, 1.

赋范空间 : IX, 3, 3.

闭路空间 : X, 3, 4.

道路空间 : X, 3, 4.

维数为 \(p\) 的射影线性簇所成的空间（在复射影空间中，维数为 \(n: \text{VIII}, 4, 3\) .

维数为 \(p\) 的射影线性簇所成的空间（在实射影空间中，维数为 \(n: \mathrm{VI}, 3, 5\) .

完备正规空间 : IX, 4, 习题 7.

波兰空间 : IX, 6, 1.

伪紧空间 : IX, 1, 习题 21.

实紧空间 : X, 4, 习题 17.

实数空间 : VI, I, I.

实射影空间 : VI, 3, 1.

可分可度量化空间 : IX, 2, 8.

苏斯林空间 : IX, 6, 2.

强零维空间 : IX, 6, 习题 1.

亚可度量化空间 : IX, 2, 习题 23.

完全非贫空间 : IX, 5, 习题 12.

超度量空间 : IX, 2, 习题 4.

零维空间 : IX, 6, 4.

球面 : IX, 2, 2, VI, 2, 3.

对数螺线 : VIII, 3, 习题 5.

正方形（闭的、开的） : VI, I, I.

\(\mathbf{R}^n\) 中的星形集 : VI, 2, 习题 12.

球极投影 : VI, 2, 4.

斯通–切赫紧化 : IX, 1, 习题 7.

斯通定理 : X, 4, 1.

斯通–魏尔斯特拉斯定理 : X, 4, 2.

严格筛分 : IX, 6, 5.

强零维空间 : IX, 6, 习题 1.

相伴子群 : VII, 1, 3.

亚可度量化空间 : IX, 2, 习题 23.

从属于子集族（函数族） : IX, 4, 3.

函数的支集 : IX, 4, 3.

角的正切 : VIII, 2, 2.

叉角的正切 : VIII, 2, 6.

数的正切 : VIII, 2, 4.

阿斯科利定理 : X, 2, 5.

贝尔定理 : IX, 5, 3.

迪尼定理 : X, 4, 1.

R. 埃利斯定理 : X, 3, 习题 25.

长田–斯米尔诺夫定理 : IX, 4, 习题 22.

斯通定理 : X, 4, 1.

斯通–魏尔斯特拉斯定理 : X, 4, 2.

乌雷松定理 : IX, 4, 1 与 4, 2.

瘦集 : IX, 5, 习题 2.

紧开拓扑 : X, 3, 4.

紧收敛的拓扑 : X, 1, 3 与 3, 4.

逐点收敛的拓扑 : X, 1, 3.

在 X 的子集上的逐点收敛的拓扑 : X, 1, 3.

S-收敛的拓扑 : X, 1, 2.

一致收敛的拓扑 : X, 1, 1.

在 X 的子集上的一致收敛的拓扑 : X, 1, 3.

一维环面 : V, 1, 2.

n 维环面 : VII, 1, 4.

旋转环面 : VII, 1, 习题 15.

完全非贫空间 : IX, 5, 习题 12.

正交变换 : VI, 2, 2.

三角不等式 : IX, 1, 1.

三角不等式 : VI, 2, 1.

复数的三角形式 : VIII, 2, 2.

\(m\) 个变量的三角多项式: \(\mathbf{X}\) , 4, 4.

超度量空间 : IX, 2, 习题 4.

连续实值函数用属于给定集合的连续

函数作的一致逼近 : X, 4, 1.

在 X 的子集上的一致收敛 : X, 1, 3.

在 \(\mathfrak{S}:\mathbf{X},\mathbf{1},\mathbf{2}\) 的各集合上的一致收敛

一致收敛点 : X, 1, 习题 9.

一致收敛的拓扑: X, 1, 1.

一致收敛的一致结构: X, 1, 1.

加法一致结构（ \(\mathbf{R}^n\) 的）: VI, 1, 2.

有界收敛的一致结构 : X, 1, 3.

紧收敛的一致结构 : X, 1, 3.

由伪度量族定义的一致结构 : IX, 1, 2.

由伪度量定义的一致结构 : IX, 1, 2.

有限开覆盖的一致结构 : IX, 4, 习题 17.

逐点收敛的一致结构 : X, 1, 3.

在 X 的子集上的逐点收敛的一致结构 : X, 1, 3.

预紧收敛的一致结构 : X, 1, 3.

\(\mathfrak{S}\) -收敛的一致结构 : X, 1, 2.

一致收敛的一致结构 : X, I, I.

在 X 的子集上的一致收敛的一致结构 : X, 1, 3.

在 S 的各集合上的一致收敛的一致结构 : X, 1, 2.

万有一致结构 : IX, 1, 习题 5.

一致收敛的（滤子） : X, I, I.

在 \(\mathfrak{S}\) 的每个集合上一致收敛的（映射族、滤子、序列、级数）: X, 1, 2.

一致等度连续的映射集 : X, 2, 1.

一致可和族 : X, 1, 2.

单位球、单位球面 : VI, 2, 3 与 IX, 3, 3.

角度量的单位 : VIII, 2, 3.

万有一致结构 : IX, 1, 习题 5.

乌雷松定理: IX, 4, 1 与 4, 2.

实赋值 : IX, 3, 2.

绝对值 : IX, 3, 2.

绝对值（复数的）: VIII, 1, 1.

半绝对值 : IX, 3, 习题 9.

赋值除环 : IX, 3, 2.

坐标簇 : VI, 1, 4.

方向向量（直线或射线的）: VI, 1, 4.

正交向量子空间 : VI, 2, 2.

正交向量 : VI, 2, 2.

球极投影的顶点 : VI, 2, 4.

魏尔斯特拉斯–斯通定理 : X, 4, 2.

零维空间 : IX, 6, 4
